% ============================================================
% Methodology appendix/section -- Two-dimensional learned warm starts
% Thesis: Frequency Helmholtz Transfer Operator
% Author: F.E. Kiewiet de Jonge
% ============================================================
%
% Suggested packages if this file is used standalone:
%   \usepackage{amsmath, amssymb}
%   \usepackage{booktabs}
%   \usepackage{graphicx}
%   \usepackage{xcolor}
%
% This file is written as a thesis-ready methodology section. It extends the
% corrected one-dimensional spectral methodology to the two-dimensional
% production problem.

\section{Two-dimensional warm-start methodology}
\label{sec:methodology_2d_warm_start}

The purpose of the two-dimensional experiments is to test whether information
learned from a lower-frequency Helmholtz solve can reduce the cost of a
high-frequency Krylov solve. The central object is a learned transfer operator
\begin{equation}
    T_{\mathrm{up}}:
    u_{\omega_L} \mapsto x_0 \approx u_{\omega_H},
    \qquad \omega_H > \omega_L,
    \label{eq:twod_transfer_operator}
\end{equation}
where $x_0$ is used as the initial guess for GMRES or flexible GMRES applied
to the high-frequency system
\begin{equation}
    A_H u_H = f,
    \qquad A_H = A(\omega_H).
    \label{eq:twod_high_frequency_system}
\end{equation}
The scientific claim is not that the neural network produces a visually
accurate field. The claim is stronger and solver-native: a useful warm start
should reduce the initial error and residual in directions that matter for
Krylov convergence, thereby reducing the number of iterations required to
reach a prescribed tolerance.

The two-dimensional setting is the decisive test. In one dimension the
matrices are small and complete eigendecomposition is possible, which makes
the problem ideal for validating the spectral interpretation. In two
dimensions the matrix on a $512 \times 512$ grid has $262{,}144$ unknowns, so
full eigendecomposition is infeasible. However, the two-dimensional GMRES
iteration count is meaningful: the system is large, indefinite, non-normal due
to the PML, and preconditioned by a realistic complex shifted Laplacian. The
headline result of the thesis must therefore be a two-dimensional solver
improvement, not only a one-dimensional spectral plot.


% ============================================================
\subsection{Two-dimensional computational domain}
\label{subsec:methodology_2d_grid}
% ============================================================

The computational grid has
\[
    n_x = n_y = 512
\]
points, so the full system contains
\[
    N = n_x n_y = 262{,}144
\]
complex unknowns. The Perfectly Matched Layer (PML) occupies
\[
    n_{\mathrm{pml}} = 112
\]
grid points on every side. The physical interior therefore has size
\[
    288 \times 288 = 82{,}944
\]
grid points. Source locations are sampled inside the physical interior, while
the full PML matrix is used when solving the PDE and when evaluating GMRES.

The model is trained and evaluated on frequency transfers
\[
    16 \to 32,
    \qquad
    32 \to 64,
    \qquad
    64 \to 128.
\]
The pair $32 \to 64$ is treated as the main diagnostic pair because previous
experiments identified it as a difficult transfer case. A method that improves
this pair is more likely to be solving a genuine bottleneck rather than only
benefiting from an easy low-frequency setting.


% ============================================================
\subsection{Role of the one-dimensional study}
\label{subsec:methodology_2d_role_of_1d}
% ============================================================

The one-dimensional experiments determine which mechanisms are worth testing
in two dimensions. Four axes of variation are especially important:
\begin{table}[h]
\centering
\begin{tabular}{lll}
\toprule
Axis & Comparison & Question \\
\midrule
Data distribution & Green function vs. FD/PML & Does PML-consistent data matter? \\
Loss masking & Full grid vs. interior only & Should the network learn PML values? \\
Input masking & Raw input vs. zeroed PML input & Does PML input contamination hurt? \\
Output zeroing & Raw output vs. zeroed PML output & Does boundary output pollute GMRES? \\
\bottomrule
\end{tabular}
\caption{Mechanisms diagnosed in one dimension before selecting the
two-dimensional training strategy.}
\label{tab:twod_axes_from_1d}
\end{table}

The resulting decision rule is simple. If a Green-function model with output
PML zeroing performs comparably to a model trained on FD/PML data, then the
existing two-dimensional Green-function dataset may be sufficient. This is
important because regenerating large two-dimensional FD/PML datasets is
expensive. If, however, the FD/PML-trained model wins clearly in initial
residual and GMRES iteration count, then the two-dimensional training data
should be regenerated with the FD/PML solver.

The one-dimensional spectral component weighting is not transferred directly
to two dimensions. In one dimension, Dirichlet eigenvectors form a clean
orthonormal basis and can be used to decompose the initial error. In two
dimensions, the full $262{,}144 \times 262{,}144$ eigendecomposition is
infeasible. The role of the one-dimensional study is therefore explanatory:
it validates the spectral mechanism and guides which warm-start variants are
worth testing in the production-scale two-dimensional solver.


% ============================================================
\subsection{PML training and evaluation approaches}
\label{subsec:methodology_2d_pml_approaches}
% ============================================================

The PML creates a methodological choice. The full numerical solver includes
the PML, but the physical quantity of interest is the interior wave field.
Training on the full grid may help the network match the exact numerical
solution, but it can also spend capacity learning boundary-layer decay rather
than interior phase and amplitude. Conversely, ignoring the PML may give a
cleaner physical objective but can create a mismatch between the warm start
and the matrix used by GMRES.

The candidate approaches are summarized in
Table~\ref{tab:twod_pml_approaches}. They are kept conceptually aligned with
the one-dimensional ablations.
\begin{table}[h]
\centering
\begin{tabular}{llll}
\toprule
Label & Training data & PML treatment & Purpose \\
\midrule
A raw &
Green function &
Keep PML input and output &
Negative control for boundary contamination \\
B zero-PML &
Green function &
Zero output PML before GMRES &
Cheap baseline using existing data \\
C PML-trained &
FD/PML &
Full-grid loss; keep output &
Tests whether learning the numerical PML helps GMRES \\
D masked &
Green function &
Zero PML input and output &
Tests pure interior-to-interior learning \\
E PML-interior &
FD/PML &
Interior loss; zero output PML &
Main physically clean PML-consistent approach \\
\bottomrule
\end{tabular}
\caption{PML approaches for the two-dimensional warm-start study.}
\label{tab:twod_pml_approaches}
\end{table}

The advantages and disadvantages are as follows.
\begin{table}[h]
\centering
\begin{tabular}{lll}
\toprule
Approach & Advantages & Disadvantages \\
\midrule
A raw &
No extra processing; tests current model directly &
Can inject large artificial values into PML rows \\
B zero-PML &
Cheap; uses existing Green-function data; removes boundary output &
Training distribution still differs from FD/PML benchmark input \\
C PML-trained &
Matches the full numerical solver; may reduce residual most &
May waste capacity on boundary-layer structure \\
D masked &
Cleanest interior-only Green-function ablation &
May discard useful incoming low-frequency context near boundary \\
E PML-interior &
PML-consistent data with physical interior objective &
Requires expensive FD/PML data generation \\
\bottomrule
\end{tabular}
\caption{Pros and cons of the PML handling strategies.}
\label{tab:twod_pml_pros_cons}
\end{table}

The main recommended two-dimensional strategy is approach E if the
one-dimensional corrected experiments show that FD/PML data distribution
matters. If approach B matches E or C in solver metrics, then B is preferred
for efficiency because it avoids regenerating the full two-dimensional
dataset.


% ============================================================
\subsection{Spectral analysis in two dimensions}
\label{subsec:methodology_2d_spectral_analysis}
% ============================================================

Complete spectral decomposition is not feasible in two dimensions. The
spectral methodology therefore changes from complete component weighting to a
combination of scalable spectral sampling and solver-native diagnostics.

\paragraph{Full eigendecomposition is avoided.}
The full matrix has $262{,}144$ unknowns. Dense eigendecomposition would scale
cubically and require memory far beyond the available resources. Moreover,
the PML matrix is complex and non-normal, so its right eigenvectors would not
provide a clean orthonormal basis for coefficient weighting even if they could
be computed.

\paragraph{Spectrum sampling is used as a diagnostic.}
Instead of all eigenvalues, sparse eigensolvers can estimate a small number of
extreme eigenvalues, for example the $100$ smallest and $100$ largest
eigenvalues or Ritz values of selected operators. This is sufficient to check
whether the broad real-part structure agrees with the validated
one-dimensional expectation. These sampled eigenvalues are diagnostic only;
they are not used for modal coefficient claims.

\paragraph{Solver-native metrics replace modal coefficients.}
The primary two-dimensional evidence is based on quantities that can be
computed for every test problem:
\begin{align}
    \text{initial residual ratio}
    &=
    \frac{\|A_H x_0 - f\|_2}{\|f\|_2},
    \label{eq:twod_initial_residual}
    \\
    \text{interior relative error}
    &=
    \frac{\|x_0[\Omega_{\mathrm{int}}]
            - u_H[\Omega_{\mathrm{int}}]\|_2}
         {\|u_H[\Omega_{\mathrm{int}}]\|_2},
    \label{eq:twod_initial_error}
    \\
    \text{GMRES iteration count}
    &=
    \min \{m : \|r_m\|_2/\|r_0\|_2 \leq \tau\}.
    \label{eq:twod_iteration_count}
\end{align}
These metrics directly measure what the thesis claims: whether the warm start
reduces the cost of solving the high-frequency linear system.

\paragraph{Residual curves are mandatory.}
A single iteration count can hide important behavior such as an initially bad
phase that GMRES corrects quickly. Therefore every serious comparison reports
the full residual curve, ideally with medians and interquartile bands across
multiple source fields.


% ============================================================
\subsection{CSL preconditioning in two dimensions}
\label{subsec:methodology_2d_csl}
% ============================================================

The baseline iterative method uses the complex shifted Laplacian (CSL)
preconditioner
\begin{equation}
    M_{\mathrm{CSL}}
    =
    A_H - i\beta\omega_H^2 I.
    \label{eq:twod_csl}
\end{equation}
GMRES is then applied to the right- or left-preconditioned system, depending
on the implementation. The shift parameter $\beta$ controls the strength of
the preconditioner. A larger $\beta$ typically gives a more damped and easier
auxiliary problem, while a smaller $\beta$ makes the preconditioner closer to
the original Helmholtz operator but harder to invert.

The warm start and CSL preconditioner have different roles:
\begin{table}[h]
\centering
\begin{tabular}{lll}
\toprule
Object & Acts on & Main effect \\
\midrule
Warm start $x_0$ & Initial guess & Reduces initial error and residual \\
CSL $M_{\mathrm{CSL}}$ & Krylov residuals & Changes convergence dynamics \\
\bottomrule
\end{tabular}
\caption{Separation between learned warm starts and classical CSL
preconditioning.}
\label{tab:twod_warmstart_vs_csl}
\end{table}

Sensitivity in $\beta$ is important. If $\beta$ is very strong, GMRES may
converge quickly regardless of the initial guess, hiding the warm-start
effect. If $\beta$ is too weak, all methods may converge slowly. Therefore
the evaluation should include at least strong, medium, and weak CSL settings,
for example $\beta \in \{0.3,0.1,0.05\}$.


% ============================================================
\subsection{Families of tests and phase structure}
\label{subsec:methodology_2d_families_phases}
% ============================================================

The experimental programme is organized into three phases. The purpose is to
avoid confusing neural-network validation loss with actual solver benefit.

\subsubsection{Phase 1: Bottleneck map}

Phase 1 diagnoses what currently limits the warm-start operator. It asks
whether performance is limited by training time, model capacity, architecture,
or objective mismatch. The main families are:
\begin{table}[h]
\centering
\begin{tabular}{lll}
\toprule
Family & Question & Examples \\
\midrule
A: capacity and budget &
Is the model undertrained or too small? &
More epochs, wider base channels, deeper UNet \\
B: operator-weighted loss &
Does solver-aligned loss help? &
Add residual-like penalty with several weights \\
Diagnostics &
Are scales and datasets sane? &
RMS, field magnitudes, source statistics \\
Evaluation &
Does field improvement reduce FGMRES cost? &
Initial residual, iterations, residual curves \\
\bottomrule
\end{tabular}
\caption{Phase 1 bottleneck families.}
\label{tab:twod_phase1_families}
\end{table}

The output of Phase 1 is a table with, for every variant and frequency pair,
the best field loss, final field loss, initial residual ratio, FGMRES
iterations, and a short conclusion. The rule is that field loss alone is never
sufficient for a claim.

\subsubsection{Phase 2: Best warm-start operator}

Phase 2 turns the Phase 1 diagnosis into the strongest defensible warm-start
model. It starts on the hard pair $32\to64$ and then repeats the best variants
on all frequency pairs. The main tracks are:
\begin{table}[h]
\centering
\begin{tabular}{lll}
\toprule
Track & Question & Candidate experiments \\
\midrule
Source conditioning &
Is $u_{\omega_L}$ enough information? &
$u_L$, $u_L+f$, $f$ only \\
Exact residual loss &
Does true PDE residual improve solver behavior? &
$L_{\mathrm{field}} + \lambda\|A_H \hat u - f\|^2$ \\
Fourier coordinates &
Does spectral coordinate bias improve phase transfer? &
Add $\sin$/$\cos$ features for low modes \\
Architecture consolidation &
Which Phase 1 model is strongest? &
Best width/depth/loss combination \\
Cross-pair validation &
Does the improvement generalize? &
Repeat on $16\to32$, $32\to64$, $64\to128$ \\
\bottomrule
\end{tabular}
\caption{Phase 2 tracks for building the best warm-start operator.}
\label{tab:twod_phase2_tracks}
\end{table}

The success criterion for Phase 2 is solver-native: the chosen model must
reduce the initial residual and the FGMRES iteration count, not merely the
validation field loss.

\subsubsection{Phase 3: Learned preconditioner}

Phase 3 moves beyond a one-time initial guess. The learning problem becomes
\begin{equation}
    r_k \mapsto z_k \approx A_H^{-1} r_k,
    \label{eq:twod_learned_preconditioner}
\end{equation}
where $r_k$ is a Krylov residual and $z_k$ is a correction. This is a learned
inverse-action or learned-preconditioner problem, not a frequency-transfer
problem. Candidate residual families include random smooth residuals,
physical source-distribution residuals, actual FGMRES residuals, and residuals
augmented with source or frequency context.

The evaluation ladder is deliberately conservative:
\begin{enumerate}
    \item test a one-shot correction
    \[
        x_1 = x_0 + M_\theta(r_0);
    \]
    \item test repeated corrections outside GMRES;
    \item integrate the learned correction inside or alongside FGMRES;
    \item compare zero start, warm start only, correction only, and warm start
    plus correction.
\end{enumerate}
This phase is more ambitious and should only follow once the warm-start
baseline is understood.


% ============================================================
\subsection{Chained frequency transfer}
\label{subsec:methodology_2d_chaining}
% ============================================================

The standard transfer experiments use one hop, for example $32\to64$.
However, in practical multilevel settings one may chain transfers:
\begin{equation}
    u_{16}
    \xrightarrow{T_{\mathrm{up}}}
    \widehat u_{32}
    \xrightarrow{T_{\mathrm{up}}}
    \widehat u_{64}
    \xrightarrow{T_{\mathrm{up}}}
    \widehat u_{128}.
    \label{eq:twod_chained_transfer}
\end{equation}
Chaining introduces a different failure mode: small phase errors can compound
across transfers. The two-dimensional evaluation should therefore compare
direct and chained warm starts where possible. For example, a warm start for
$\omega=128$ can be produced either from a direct low-frequency input or by
successive applications of pairwise transfer models. The comparison should be
judged by initial residual and FGMRES iterations.


% ============================================================
\subsection{Sensitivity analysis}
\label{subsec:methodology_2d_sensitivity}
% ============================================================

Sensitivity analysis is required because a warm-start benefit can be hidden or
amplified by numerical choices unrelated to the neural operator. The key
parameters are:
\begin{table}[h]
\centering
\begin{tabular}{lll}
\toprule
Parameter & Values or range & Reason \\
\midrule
CSL shift $\beta$ & $0.3$, $0.1$, $0.05$ & Tests strong to weak preconditioning \\
PML damping scale & Around default values & Tests boundary absorption strength \\
PML treatment & A--E & Tests data, loss, input, and output choices \\
Frequency pair & All three pairs & Tests transfer difficulty and generalization \\
Source family & Single and multi-source & Tests robustness to forcing distribution \\
Architecture & Best Phase 1 variants & Separates model capacity from physics choices \\
\bottomrule
\end{tabular}
\caption{Sensitivity dimensions for the two-dimensional warm-start study.}
\label{tab:twod_sensitivity}
\end{table}

The sensitivity analysis should be staged. First, run the corrected
one-dimensional experiment to decide whether FD/PML training data are
necessary. Second, port the smallest winning set of variants to two
dimensions. Third, vary CSL strength and frequency pair. Only after these
results are understood should broader architecture and data-generation sweeps
be launched.


% ============================================================
\subsection{Decision logic for porting one-dimensional results to two dimensions}
\label{subsec:methodology_2d_decision_logic}
% ============================================================

The one-dimensional results determine the two-dimensional plan according to
the following rules:
\begin{table}[h]
\centering
\begin{tabular}{lll}
\toprule
One-dimensional observation & Interpretation & Two-dimensional action \\
\midrule
B matches C/E &
Green data plus output zeroing is enough &
Use existing 2D data first \\
C wins clearly &
Full PML prediction helps solver residual &
Generate 2D FD/PML data and train C \\
E wins or matches C &
Interior objective is sufficient and cleaner &
Train 2D FD/PML with interior loss \\
A has good field loss but bad residual &
PML contamination matters &
Never use raw PML output as final method \\
Higher pairs degrade strongly &
Transfer gap is limiting &
Prioritize $32\to64$ and source conditioning \\
\bottomrule
\end{tabular}
\caption{Decision logic for transferring conclusions from one dimension to
two dimensions.}
\label{tab:twod_decision_logic}
\end{table}

This decision table prevents unnecessary expensive data generation. The
two-dimensional FD/PML dataset should only be regenerated if the one-dimensional
evidence shows that matching the PML data distribution improves solver-native
metrics.


% ============================================================
\subsection{Reporting requirements}
\label{subsec:methodology_2d_reporting}
% ============================================================

Every two-dimensional result should report four quantities:
\begin{enumerate}
    \item initial relative residual $\|A_H x_0 - f\|_2/\|f\|_2$;
    \item interior relative error
    $\|x_0[\Omega_{\mathrm{int}}]-u_H[\Omega_{\mathrm{int}}]\|_2/
      \|u_H[\Omega_{\mathrm{int}}]\|_2$;
    \item FGMRES iterations to tolerance;
    \item residual convergence curves over iterations.
\end{enumerate}
When available, wall-clock time should also be reported, including neural
network inference cost. A method that reduces iterations but increases total
time may still be scientifically interesting, but the distinction must be
explicit.

The guiding rule is:
\[
    \text{do not trust field loss alone.}
\]
Every major claim must ultimately be supported by initial residuals, FGMRES
iteration counts, and residual curves.


% ============================================================
\subsection{Summary of the full research pipeline}
\label{subsec:methodology_2d_pipeline_summary}
% ============================================================

The full research pipeline can be summarized as follows.
\begin{enumerate}
    \item Validate the spectral mechanism in one dimension using Dirichlet
    eigenvectors with unit norm.
    \item Use the one-dimensional PML ablations to decide which data and loss
    choices are worth porting to two dimensions.
    \item In two dimensions, evaluate warm starts with solver-native metrics:
    initial residual, interior error, FGMRES iterations, and residual curves.
    \item Use Phase 1 to identify bottlenecks in the current warm-start model.
    \item Use Phase 2 to build the strongest physics-aware warm-start operator.
    \item Use Phase 3, if time permits, to move from warm starts toward learned
    inverse action on Krylov residuals.
\end{enumerate}
The clean thesis arc is therefore:
\[
    \text{diagnose the bottleneck}
    \;\longrightarrow\;
    \text{build a solver-reducing warm start}
    \;\longrightarrow\;
    \text{extend toward learned preconditioning}.
\]

